\documentclass[12pt, a4paper]{article}

% 引入必要的宏包
\usepackage[UTF8]{ctex}     % 中文支持
\usepackage{amsmath, amssymb, amsthm} % 数学符号和环境
\usepackage{geometry}       % 页面设置
\usepackage{enumitem}       % 列表定制
\usepackage{graphicx}       % 图片支持
\usepackage{fancyhdr}       % 页眉页脚

% 页面边距设置
\geometry{left=2.5cm, right=2.5cm, top=2.5cm, bottom=2.5cm}

% 宏定义：方便排版选择题选项
% 使用方法：\choices{选项A}{选项B}{选项C}{选项D}
\newcommand{\choices}[4]{
    \begin{enumerate}[label=\Alph*., itemsep=0pt, parsep=0pt, topsep=5pt, labelsep=0.5em]
        \item #1
        \item #2
        \item #3
        \item #4
    \end{enumerate}
}
% 横向排列的选项 (2列)
\newcommand{\choicesTwo}[4]{
    \begin{enumerate}[label=\Alph*., itemsep=0pt, parsep=0pt, topsep=5pt, labelsep=0.5em]
        \begin{minipage}{0.45\linewidth} \item #1 \end{minipage}
        \begin{minipage}{0.45\linewidth} \item #2 \end{minipage}
        \begin{minipage}{0.45\linewidth} \item #3 \end{minipage}
        \begin{minipage}{0.45\linewidth} \item #4 \end{minipage}
    \end{enumerate}
}
% 横向排列的选项 (4列)
\newcommand{\choicesFour}[4]{
    \begin{enumerate}[label=\Alph*., itemsep=0pt, parsep=0pt, topsep=5pt, labelsep=0.5em]
        \begin{minipage}{0.22\linewidth} \item #1 \end{minipage}
        \begin{minipage}{0.22\linewidth} \item #2 \end{minipage}
        \begin{minipage}{0.22\linewidth} \item #3 \end{minipage}
        \begin{minipage}{0.22\linewidth} \item #4 \end{minipage}
    \end{enumerate}
}

\title{\textbf{《高等数学》必刷习题——基础版}}
\author{}
\date{}

\begin{document}

\section*{第七章 定积分（基础）}

\begin{enumerate}
    \item 已知 $ f(x) $ 为 $ [1, +\infty] $ 上的连续函数，且 $ F(x) = \int_{1}^{x^{2}} \frac{f(t)}{t} dt, x \in (1, +\infty) $ 。则 $ F'(x) = $ \underline{\hspace{1cm}}。
    \choicesFour{$ 2f(x) $}{$ 2xf(x^{2}) $}{$ \frac{f(x^{2})}{x} $}{$ \frac{2f(x^{2})}{x} $}

    \item 定积分 $ \int_{-2}^{0}|x+1|dx $ 的值为 \underline{\hspace{1cm}}。
    \choicesFour{-2}{2}{-1}{1}

    \item 定积分 $ \int_{-1}^{1}\frac{x^{4}\tan x}{2x^{2}+\cos x}dx $ 的值为 \underline{\hspace{1cm}}。
    \choicesFour{0}{1}{-1}{2}

    \item 设 $ p = \int_{0}^{\frac{\pi}{2}} \sin^{2}x \, dx $ , $ q = \int_{0}^{\frac{\pi}{2}} \cos^{2}x \, dx $ , $ r = \int_{-\frac{\pi}{2}}^{\frac{\pi}{2}} \sin^{2}x \, dx $ , 则 \underline{\hspace{1cm}}。
    \choicesFour{$ p = q = r $}{$ p = q < r $}{$ p < q < r $}{$ p > q > r $}

    \item $ \int_{a}^{b}f'(2x)dx= $ \underline{\hspace{1cm}}。
    \choicesTwo
    {$ f(b) - f(a) $}
    {$ f(2b) - f(2a) $}
    {$ \frac{1}{2}[f(2b)-f(2a)] $}
    {$ 2[f(2b)-f(2a)] $}

    \item 设 $ \varphi'(x) = \int_{0}^{\ln x} t^{2} dt $ , 则 $ \varphi'(x) = $ \underline{\hspace{2cm}}。

    \item 函数 $ f(x) $ 满足 $ \int_{0}^{2x} f(t) dt = 1 + x^{3} $ ，则 $ f(8) = $ \underline{\hspace{2cm}}。

    \item $ \int_{0}^{\pi}\cos x \, dx= $ \underline{\hspace{2cm}}。

    \item 定积分 $ \int_{0}^{1}\frac{dx}{100+x^{2}}= $ \underline{\hspace{2cm}}。

    \item 定积分 $ \int_{1}^{2}\frac{\sqrt{x-1}}{x}dx= $ \underline{\hspace{2cm}}。

    \item 求定积分 $ \int_{1}^{5}e^{\sqrt{2x-1}}dx $。

    \item 求定积分 $ \int_{1}^{4}\frac{1+\ln x}{\sqrt{x}}dx $。

    \item 求定积分 $ \int_{0}^{\sqrt{3}} \arctan x \, dx $。

    \item $ f(x)=\begin{cases}xe^{-x}, & x\leq0\\ 3x^{2}, & 0<x\leq1\end{cases} $ ，求 $ \int_{-3}^{1}f(x)dx $。

    \item 求定积分 $ \int_{0}^{\pi} x(\sin x)^{3}dx $。
\end{enumerate}


\end{document}